\chapter{Preliminary Logic: Classical and Paraconsistent Systems}
\label{ch:preliminary-logic}

\section{Introduction}

The introduction of an indeterminate element $\I$ that satisfies $\I = 0/0$ challenges fundamental assumptions in classical logic and mathematics. Before we can rigorously develop the axiomatic system that accommodates $\I$, we must establish the logical framework within which our reasoning will operate.

This chapter examines two contrasting logical paradigms:
\begin{enumerate}
    \item \textbf{Classical Logic}: The traditional two-valued logic where every proposition is either true or false, governed by the principle of non-contradiction.
    \item \textbf{Paraconsistent Logic}: Systems that allow for the presence of contradictions without trivializing the entire theory (i.e., without allowing any proposition to be proven).
\end{enumerate}

The need for paraconsistent considerations arises from the fact that $\I$ appears to embody contradictory properties: it is neither zero nor any determinate value, yet it interacts with standard mathematical operations.

\section{Classical Logic Review}

\subsection{Propositional Calculus}

Classical propositional logic is based on the following connectives and their truth tables:

\begin{definition}[Logical Connectives]
Let $P$ and $Q$ be propositions. The standard connectives are:
\begin{itemize}
    \item \textbf{Negation} ($\neg P$): True when $P$ is false
    \item \textbf{Conjunction} ($P \land Q$): True when both $P$ and $Q$ are true
    \item \textbf{Disjunction} ($P \lor Q$): True when at least one of $P$ or $Q$ is true
    \item \textbf{Implication} ($P \rightarrow Q$): False only when $P$ is true and $Q$ is false
    \item \textbf{Equivalence} ($P \leftrightarrow Q$): True when $P$ and $Q$ have the same truth value
\end{itemize}
\end{definition}

\subsection{Fundamental Principles}

\begin{axiom}[Law of Non-Contradiction]
For any proposition $P$:
\[
\neg(P \land \neg P)
\]
It is impossible for a proposition and its negation to both be true.
\end{axiom}

\begin{axiom}[Law of Excluded Middle]
For any proposition $P$:
\[
P \lor \neg P
\]
Every proposition is either true or false; there is no third option.
\end{axiom}

\begin{axiom}[Principle of Explosion (Ex Falso Quodlibet)]
From a contradiction, anything follows:
\[
(P \land \neg P) \rightarrow Q
\]
for any proposition $Q$.
\end{axiom}

\subsection{Limitations for Indeterminate Elements}

The classical framework poses immediate challenges when dealing with $\I$:

\begin{example}[Division by Zero Paradox]
Consider the statement: ``$\I = 0/0$''. In classical arithmetic:
\begin{align*}
    \I &= \frac{0}{0} \\
    0 \cdot \I &= 0 \\
    0 \cdot \I &= 0 \cdot x \quad \text{for any } x
\end{align*}
This would imply $\I = x$ for all $x$, a contradiction that would trigger explosion in classical logic.
\end{example}

\section{Paraconsistent Logic}

\subsection{Motivation and Historical Development}

Paraconsistent logic emerged from the work of Newton da Costa, Graham Priest, and others who recognized that certain domains naturally involve contradictory information without collapsing into triviality.

\begin{definition}[Paraconsistent Logic]
A logic is \textbf{paraconsistent} if it rejects the principle of explosion, i.e., it is possible to have $P \land \neg P$ without deriving all propositions.
\end{definition}

\subsection{Da Costa's System $C_1$}

One of the most studied paraconsistent systems is da Costa's hierarchy $C_n$. We focus on $C_1$:

\begin{definition}[Well-Behaved Formula]
A formula $P$ is \textbf{well-behaved} (denoted $P^\circ$) if:
\[
P^\circ \equiv \neg(P \land \neg P)
\]
\end{definition}

In $C_1$, the principle of explosion holds only for well-behaved formulas:
\[
(P \land \neg P \land P^\circ) \rightarrow Q
\]

\subsection{Relevance Logic}

Another approach, relevance logic, requires that in a valid inference $P \rightarrow Q$, the premise $P$ must be actually used (relevant) in deriving $Q$.

\begin{axiom}[Relevance Condition]
In an inference $P \vdash Q$, there must be a demonstrable connection between $P$ and $Q$.
\end{axiom}

This prevents arbitrary conclusions from contradictions.

\section{The Logic of Indeterminacy}

\subsection{Three-Valued Logic for $\I$}

We propose a three-valued logic where propositions can be:
\begin{enumerate}
    \item \textbf{True} (T)
    \item \textbf{False} (F)
    \item \textbf{Indeterminate} ($\I$)
\end{enumerate}

\begin{definition}[Truth Tables for $\I$-Logic]
The extended truth tables are:

\begin{center}
\begin{tabular}{c|c|c}
$P$ & $\neg P$ & $P^\circ$ \\
\hline
T & F & T \\
F & T & T \\
$\I$ & $\I$ & F
\end{tabular}
\quad
\begin{tabular}{c|c|c|c|c}
$P$ & $Q$ & $P \land Q$ & $P \lor Q$ & $P \rightarrow Q$ \\
\hline
T & T & T & T & T \\
T & F & F & T & F \\
T & $\I$ & $\I$ & T & $\I$ \\
F & T & F & T & T \\
F & F & F & F & T \\
F & $\I$ & F & $\I$ & T \\
$\I$ & T & $\I$ & T & T \\
$\I$ & F & F & $\I$ & $\I$ \\
$\I$ & $\I$ & $\I$ & $\I$ & $\I$
\end{tabular}
\end{center}
\end{definition}

\subsection{Non-Triviality Preservation}

\begin{theorem}[Controlled Indeterminacy]
\label{thm:controlled-indeterminacy}
In $\I$-logic, the presence of indeterminate values does not lead to trivialization. Specifically, if $P$ has truth value $\I$, we cannot derive arbitrary $Q$ with determinate truth value.
\end{theorem}

\begin{proof}
Consider the inference $(P \land \neg P) \rightarrow Q$ where $P$ has value $\I$:
\begin{align*}
    P \land \neg P &= \I \land \I = \I \\
    \I \rightarrow Q &= \begin{cases}
        T & \text{if } Q = T \\
        \I & \text{if } Q = F \\
        \I & \text{if } Q = \I
    \end{cases}
\end{align*}
The implication does not force $Q$ to be true when $Q$ is false, preventing explosion.
\end{proof}

\section{Consistency vs. Non-Triviality}

\subsection{Defining Consistency}

\begin{definition}[Consistency]
A theory $T$ is \textbf{consistent} if there exists no formula $P$ such that both $T \vdash P$ and $T \vdash \neg P$.
\end{definition}

\begin{definition}[Non-Triviality]
A theory $T$ is \textbf{non-trivial} if there exists at least one formula $Q$ such that $T \nvdash Q$.
\end{definition}

In classical logic, consistency implies non-triviality (by the principle of explosion). In paraconsistent logic, we can have:
\[
\text{Inconsistent but Non-Trivial}
\]

\subsection{The Role of $\I$ in Preserving Non-Triviality}

\begin{theorem}[Non-Triviality Preservation]
\label{thm:nontriviality}
The extended theory $\text{ZF} + \text{AC} + \I_1 + \cdots + \I_{15}$ is non-trivial.
\end{theorem}

\begin{proof}[Proof Sketch]
We construct a model where:
\begin{enumerate}
    \item All standard ZF + AC axioms are satisfied
    \item The element $\I$ exists and satisfies $\kappa \ominus \kappa = \I$
    \item There exist propositions that are not provable (e.g., $\I = 0$)
\end{enumerate}
The existence of such a model, which we will construct in detail in Chapter~\ref{ch:model-constructions}, guarantees non-triviality.
\end{proof}

\section{The Corollary of Triviality: Why $\I \neq 0$}

\subsection{The Danger of $\I = 0$}

\begin{theorem}[Triviality Collapse]
\label{thm:triviality-collapse}
If $\I = 0$, then the extended theory collapses to the trivial theory where all propositions are provable.
\end{theorem}

\begin{proof}
Assume $\I = 0$. From Axiom $\I_1$, we have:
\[
\kappa \ominus \kappa = \I = 0
\]
for any cardinal $\kappa$. Consider $\kappa = 1$:
\[
1 \ominus 1 = 0
\]
This is consistent with classical arithmetic. However, from Axiom $\I_4$ (Absorption):
\[
\I \oplus \kappa = \I
\]
If $\I = 0$, this becomes:
\[
0 \oplus \kappa = 0 \quad \Rightarrow \quad \kappa = 0
\]
for all cardinals $\kappa$. This forces all cardinals to equal zero, collapsing the entire cardinal structure.

In particular, $1 = 0$, which in a ring structure implies:
\[
\forall x, \quad x = x \cdot 1 = x \cdot 0 = 0
\]
Thus every element equals zero, giving us the trivial theory.
\end{proof}

\subsection{Philosophical Implications}

The requirement that $\I \neq 0$ is not merely a technical condition but a fundamental preservation of mathematical richness:

\begin{itemize}
    \item \textbf{Ontological}: $\I$ represents genuine indeterminacy, not the absence of value
    \item \textbf{Epistemological}: Our knowledge admits states of fundamental uncertainty
    \item \textbf{Mathematical}: Non-trivial structures require genuine distinction between elements
\end{itemize}

\section{Inference Rules in $\I$-Logic}

\subsection{Modified Modus Ponens}

\begin{axiom}[Guarded Modus Ponens]
\[
\frac{P, \quad P \rightarrow Q, \quad P^\circ}{Q}
\]
Modus ponens requires that $P$ be well-behaved.
\end{axiom}

\subsection{Universal Instantiation}

\begin{axiom}[Restricted Universal Instantiation]
\[
\frac{\forall x \, \phi(x), \quad x^\circ}{\phi(x)}
\]
When instantiating universal statements, the instance must be well-behaved.
\end{axiom}

\section{Conclusion}

This chapter has established the logical foundations necessary for developing a rigorous theory of the indeterminate element $\I$. Key insights include:

\begin{enumerate}
    \item Classical logic is insufficient due to the principle of explosion
    \item Paraconsistent logic provides a framework for controlled contradictions
    \item A three-valued logic naturally accommodates indeterminate truth values
    \item The requirement $\I \neq 0$ is essential for non-triviality
    \item Modified inference rules preserve soundness while allowing indeterminacy
\end{enumerate}

With these logical tools in hand, we proceed to the detailed axiomatization in the following chapters, where each of the 15 axioms will be examined individually with full mathematical rigor.

\section{Exercises}

\begin{enumerate}
    \item Prove that in $\I$-logic, $P \lor \neg P$ can have truth value $\I$.
    \item Show that if both $P$ and $P \rightarrow Q$ have truth value $\I$, we cannot determine the truth value of $Q$ without additional information.
    \item Construct a model of $C_1$ that includes an element satisfying the properties of $\I$.
    \item Verify that the truth tables given for $\I$-logic satisfy associativity for conjunction and disjunction.
    \item Prove that if $\I = 1$ (rather than $\I = 0$), a different form of trivialization occurs.
\end{enumerate}
